\subsection{Potenziali separabili in coordinate cartesiane}
Il problema $ n $-dimensionale si riduce a $ n $ problemi 1-dimensionali.
\[
V(\vec{x}) = \sum_i V_i(x^{(i)})
\]
N.B.: $ \hat{T} $ \`e sempre separabile.
\[
\mathcal{H} = \frac{{\hat{\vec{p}}}^2}{2m} + \sum_i V_i(\hat{x}^{(i)}) = \sum_i \mathcal{H}_i,
\]
\[
\mathcal{H}_i = \frac{{\hat{p}^{(i)}}^2}{2m} + V_i(\hat{x}^{(i)}).
\]

Allora faccio
\begin{align*}
\braket{\vec{x}|\mathcal{H}|\psi} &= \braket{\vec{x}|\mathcal{H}_1|\psi} + \braket{\vec{x}|\mathcal{H}_2|\psi} + \cdots = \\
&= \left( -\frac{\hbar^2}{2m} \frac{d^2}{d{x^{(1)}}^2} + V(\hat{x}_1) \right) \psi(\vec{x}) \\
&+ \left( -\frac{\hbar^2}{2m} \frac{d^2}{d{x^{(2)}}^2} + V(\hat{x}_2) \right) \psi(\vec{x}) + \cdots
\end{align*}

Se inoltre ho $ \mathcal{H} \ket{\psi} = E \ket{\psi} $,
\[
\mathcal{H} \ket{\psi_n} = E_n \ket{\psi_n}, \quad \braket{x^{(i)}|\psi_n} = \psi_n(x^{(i)}).
\]
Se \`e $ d $ il numero di dimensioni, allora 
\begin{align*}
\psi_{n_1, \dots, n_d}(\vec{x}) &= \psi_{n_1}(x^{(1)}) \cdots \psi_{n_d}(x^{(d)}) \\
\braket{\vec{x}|\mathcal{H}|\psi_{n_1, \dots, n_d}}  &= \frac{\hbar^2}{2m} \left( \frac{d^2}{d{x^{(1)}}^2} \psi_{n_1}(x^{(1)}) \right) \psi_{n_2}(x^{(2)}) \cdots \psi_{n_d}(x^{(d)}) + \\
&+ V_1(x^{(1)}) \psi_{n_1}(x^{(1)}) \cdots \psi_{n_d}(x^{(d)}) + \\
&- \frac{\hbar^2}{2m} \psi_{n_1}(x^{(1)}) \left( \frac{d^2}{d{x^{(2)}}^2} \psi_{n_2}(x^{(2)}) \right) \psi_{n_3}(x^{(3)}) \cdots \psi_{n_d}(x^{(d)}) + \cdots \\
&= (E_{n_1} + \cdots + E_{n_d}) \psi_{n_1, n_2, \dots, n_d}(\vec{x}) \quad \Rightarrow
\end{align*}
\[
\Rightarrow \quad -\frac{\hbar^2}{2m} \frac{d^2}{d{x^{(i)}}^2} \psi_{n_i}(x^{(i)}) + V_i(x^{(i)})\psi_{n_i}(x^{(i)}) = E_{n_i} \psi_{n_i}(x^{(i)})
\]